\chapter{Informações Gerais do Trabalho}

As informações gerais, também chamadas de metadados, ficam no preâmbulo de \texttt{tcc.tex}. Elas alimentam páginas iniciais e propriedades do PDF.

\section{Autoria, título e vínculo}

\begin{lstlisting}
\autor{Nome do Aluno}
% \autores{Aluno 1//Aluno 2}
\titulo{Título do Trabalho}
\subtitulo{Subtítulo opcional}
\curso{Nome do curso}
\tipotrabalho{Trabalho de Conclusão de Curso}
\graupretendido{bacharel em Nome do curso}
\local{Cidade}
\ano{Ano}
\end{lstlisting}

Use \texttt{\textbackslash autor} para uma autoria. Para múltiplas autorias admitidas pelo curso, use \texttt{\textbackslash autores} e separe nomes com \texttt{//}. Não use os dois comandos ao mesmo tempo.

\section{Orientação e aprovação}

\begin{lstlisting}
\orientador{Nome do Orientador}
\coorientador{}
\areadeconcentracao{Área de concentração}
\dataaprovacao{Brasília, DF, dia de mês de ano}
\end{lstlisting}

Deixe o argumento de \texttt{\textbackslash coorientador} vazio quando não houver coorientação. A data e o formato da banca devem acompanhar a decisão institucional aplicável; os valores do modelo são exemplos editáveis.

\section{Banca examinadora}

Limpe a lista padrão e cadastre cada membro com nome, prefixo acadêmico e vínculo:

\begin{lstlisting}
\limparmembrosbanca
\membrobanca{Nome do Orientador}{Prof. Dr.}{UnDF}
\membrobanca{Nome do Membro 1}{Profa. Dra.}{Instituição}
\membrobanca{Nome do Membro 2}{Prof. Dr.}{Instituição}
\end{lstlisting}

O primeiro argumento contém o nome sem titulação. O segundo é impresso antes do nome. O terceiro aparece na linha de vínculo institucional. Cadastre apenas a composição efetivamente aprovada.

\section{Conferência}

Depois de preencher os campos, confira capa, folha de rosto, folha de aprovação e propriedades do PDF. Um erro no preâmbulo pode aparecer em mais de uma página.
